\documentclass[10pt,a4paper]{article}
\usepackage[utf8]{inputenc}
\usepackage[T1]{fontenc}
\usepackage{lmodern}
\usepackage[margin=2.3cm]{geometry}
\usepackage{amsmath}
\usepackage{graphicx}
\usepackage{booktabs}
\usepackage{caption}
\usepackage{float}
\usepackage{hyperref}
\usepackage{parskip}
\captionsetup{font=small,labelfont=bf}
\input{values}

\title{From transition lines to charge states:\\
what the U-Net line maps are good for\\[4pt]
\large the question, what was done, what came out, what it means}
\author{\small Generated by \texttt{docs/charge\_states/make\_report.py}\\
\small run: \texttt{\vRun}}
\date{}

\begin{document}
\maketitle

\section*{Summary}

\textbf{The question.}  The U-Net predicts the transition lines of a double
quantum dot from a few percent of the gate plane.  But a line map is not
what an experimentalist needs.  They need to know \emph{which charge state
$(n_1,n_2)$ the device is in at each gate voltage}, and in particular where
the $(1,1)$ state is, so that the gates can be set there.  The question
raised by Prof.\ Fujita is what predicting the lines directly gains, or
loses, compared with the conventional route, and whether the result is
useful for the next step of the analysis.

\textbf{What was done.}  A program was written that reads the charge state of
every pixel off a line map, using only the lines.  It was applied to the
U-Net maps and scored against the simulator's exact charge states.  Two
references were added: the conventional route (measure the whole image, find
the lines without a network) and a trivial floor (call everything $(0,0)$).
Because the U-Net maps fail through gaps, a geometric repair step was added
that fits straight lines and closes the gaps.  The evaluation focuses on
the few-electron region ($n_1,n_2\le\vNmax$) and on two numbers an
experiment cares about: how many charge states are identified correctly
(balanced accuracy), and whether the gates, set at the predicted $(1,1)$,
really land in $(1,1)$.

\textbf{What came out.}  For the largest budget ($8\times60$ rays $\times$
points, \vEqNPts{} measurements, \vUCov\,\% of the grid):
\begin{itemize}
\item U-Net map as is: balanced accuracy \vUBal\,\%, landing \vULand\,\%;
\item U-Net map with the line fit: \vFBal\,\% and \vFLand\,\%;
\item conventional route with the \emph{same} \vEqNPts{} measurements, best
  of the variants tried: \vEqCompBal\,\%;
\item conventional route on the full image (\vFold{} times more
  measurements): \vDBal\,\% and \vDLand\,\%.
\end{itemize}
At equal measurement cost, the U-Net with the line fit beats every
conventional variant tried, at all \vEqNBud{} budgets, by \vEqDiff{} points
at $8\times60$ (95\,\% interval \vEqDiffLo--\vEqDiffHi).  Measuring the whole
image is still far better.

\textbf{What it means.}  This is the first evidence that predicting lines
directly is worth something: from few measurements it recovers the charge
states much better than reconstructing an image from the same measurements
and then detecting lines.  Two observations matter as much as the numbers.
First, the usual line-level score (F1) does \emph{not} predict whether a
U-Net map is usable for charge states; what decides it is whether the lines
are closed.  Second, the advantage depends on the line fit; the raw U-Net
map alone is better than the conventional route at only \vEqRawWins{} of
\vEqNBud{} budgets.  The competitors here are simple ones, so the claim is
not yet strong enough to publish without the checks listed in
Sec.~\ref{sec:contribution}.

\section{The question}

\subsection{Why charge states}
Tuning a double quantum dot means finding gate voltages where each dot
holds a known number of carriers, typically one each, $(1,1)$, for a
spin qubit.  The charge stability diagram (CSD) is the tool for this: its
transition lines divide the gate plane into regions of constant
$(n_1,n_2)$.  Finding the lines is a means; the end is the map of charge
states, and above all the location of $(1,1)$.

\subsection{What the U-Net does}
The project's U-Net measures the charge sensor along a few straight rays
through the gate plane (\vCovLo--\vCovHi\,\% of the $100\times100$ grid) and
predicts a probability $P$ that each pixel lies on a transition line.  With
a threshold $t$, $B=P>t$ is a binary line map.  At the line level it works:
F1@1 (a line found within one pixel) is \vLineFoneMin--\vLineFoneMax{} on the
test devices, depending on the budget.

\subsection{The question raised}
Prof.\ Fujita's point, in short: in simulation the transition lines are
known exactly, so finding them is not the result in itself.  The
conventional route, reconstructing the image and then detecting the lines,
also preserves the transitions.  What does predicting the lines directly
gain or lose, and is the result useful for what comes after?

This report answers the second half directly, by \emph{using} the line maps
for what comes after (assigning charge states), and measures how useful
they are.  The first half (gain or loss against the conventional route) is
answered only partly, and Sec.~\ref{sec:contribution} explains what is
missing.

\section{What was done, and why}

\subsection{The data and the ground truth}
The study uses the existing 4--8-ray sweep: 15 trained U-Nets (4 to 8 rays
$\times$ 40 to 60 points per ray), 500 training devices (75 of them a
validation slice) and \vNTest{} held-out test devices.  For every device the
simulator (QArray) also gives the exact occupation $(n_1,n_2)$ at every
pixel, from the capacitance model with no sensor involved.  This is the
answer key: it is used only to score, never as an input.

The devices are very different from each other.  The capacitances are drawn
at random over a wide range and the voltage window is shifted at random, so
a test window holds \vNStMin{} to \vNStMax{} charge states (median
\vNStMed).  On every device the top-right corner of the window is $(0,0)$:
the carriers are holes, so the occupation grows towards negative voltage.
This was checked on the data, not assumed.

\subsection{Reading charge states off a line map}
The program \texttt{label\_states} works on a line map alone and does what an
experimentalist does by eye (Fig.~\ref{fig:method}):
\begin{enumerate}
\item the lines cut the plane into regions; each region is one charge state;
\item every border between two regions is classified by its slope: the
  steep lines change $n_1$ (dot~1), the shallow ones $n_2$ (dot~2), and the
  short ones with positive slope move a carrier between the dots
  (interdot);
\item the side at lower voltage holds the extra hole;
\item the top-right region is $(0,0)$, and the program counts outwards
  across the borders, along the longest (most reliable) borders first.
\end{enumerate}
Every border that was not used for counting is a second route between two
regions that already have labels.  Whether it agrees is a self-check, the
\emph{consistency}, which needs no ground truth and could be used in an
experiment.

\begin{figure}[!htbp]
\centering
\includegraphics[width=\textwidth]{figures/step3_tree_labels.pdf}
\caption{Reading charge states off a line map, one test device.  Top: the
true line map.  Bottom: the U-Net map.  Left: the regions and the borders
used for counting (solid, from the yellow $(0,0)$ star) or only checked
(dashed; red where they disagree).  Middle: the labels.  Right: the truth.
On the U-Net map, gaps merged most regions into one.}
\label{fig:method}
\end{figure}

\paragraph{A check that found a bug.}  The program must be perfect on the
true line map, otherwise every later number would be partly the program's
fault.  At first it was not.  On the first training devices checked during
development the accuracy was below 40\,\%.  The cause was that the
simulator's line maps stop one pixel short of the image edge, so regions
joined around the ends of the lines.  After closing the lines onto the
frame, the true line maps are labelled essentially perfectly (100\,\% of
the pixels, 99\,\% of the states on the test devices).  From then on, any
error is the line map's, not the program's.

\subsection{The conventional route, as a competitor}
To know whether the U-Net helps, it must be compared with what one would
do without it: measure every pixel of the sensor image and find the lines
directly.  This was harder than expected.  Standard edge detection (Canny)
and gradient thresholds failed.  The sensor's own smooth background rises
steeply at the edge of the window, and these detectors draw that instead of
the transitions.  What works is to look for \emph{steps}: compare each
pixel-to-pixel difference with the local background slope.  A transition
is a jump on top of it.  Out of \vNDense{} settings of the three detectors,
the one used was chosen on the validation devices only
(Fig.~\ref{fig:dense}).

\begin{figure}[!htbp]
\centering
\includegraphics[width=0.92\textwidth]{figures/dense_examples.pdf}
\caption{The conventional route: (a) the full sensor image, (b) the step
residual, (c) the line map from it, (d) the charge states read off (c) by the
same program used for the U-Net.}
\label{fig:dense}
\end{figure}

\subsection{Why the U-Net maps fail, and a repair}
On the U-Net maps the labels were mostly wrong.  Looking at the failures
showed one dominant cause.  The U-Net finds the lines, but leaves
\emph{gaps} in them (Fig.~\ref{fig:regions}).  A single gap joins two charge
states into one region, and the error is then carried into every region
counted beyond it.

\begin{figure}[!htbp]
\centering
\includegraphics[width=\textwidth]{figures/step1_regions.pdf}
\caption{Where the failure starts.  Top: true line map and its regions.
Bottom: the U-Net probability (e) shows the lines, but after the threshold
(f) they have gaps, and the gaps merge the regions (g, h).}
\label{fig:regions}
\end{figure}

The repair, \texttt{fit\_lines}, uses only geometry (Fig.~\ref{fig:fit}):
\begin{enumerate}
\item thin the U-Net lines to one pixel and cut them into pieces at the
  crossings;
\item fit a straight line to each piece (least squares);
\item continue each piece that simply stops, along its own direction;
\item keep the continuation only if it reaches another line or the edge of
  the window.
\end{enumerate}
Its three settings were chosen on the validation devices
(\emph{\vFitKnob}, out of \vNFitSet).

\begin{figure}[!htbp]
\centering
\includegraphics[width=\textwidth]{figures/fit_steps.pdf}
\caption{The line fit, step by step, on one test device.  Red continuations
(5) close the gaps, and the labels improve from (7) to (8).  The diagonal red
continuation on the left is a false line: it cuts the $(1,0)$ region and
causes the wrong $(0,0)$ patch in (8).}
\label{fig:fit}
\end{figure}

\subsection{How it was scored}
Three things were decided on purpose, because the obvious metric is
misleading here.
\begin{itemize}
\item \textbf{The few-electron region only} ($n_1,n_2\le\vNmax$), with
  absolute labels.  An experiment works there and needs the real occupation.
\item \textbf{Balanced accuracy}, not pixel accuracy.  The $(0,0)$ region is
  large, and calling everything $(0,0)$ already gets most pixels right.
  Balanced accuracy computes the accuracy per charge state and averages
  them, so every state counts once.  ``All $(0,0)$'' then scores
  \vBaseBal\,\%.
\item \textbf{(1,1) landing}: set the gates at the centre of the predicted
  $(1,1)$ region and ask whether the device is really in $(1,1)$ there.
  This is the yes/no question of tuning.
\end{itemize}
Everything that has a setting (the U-Net threshold, the dense detector, the
line fit) was chosen on validation devices; the test devices were only used
to report.  The averages come with 95\,\% bootstrap intervals over the test
devices.  Details are in \texttt{evaluation\_method.pdf}
(Fig.~\ref{fig:eval}).

\begin{figure}[!htbp]
\centering
\includegraphics[width=\textwidth]{figures/eval_states.pdf}
\caption{Scoring one device: truth, predicted labels, right and wrong
pixels, the landing point, and the per-state accuracies whose mean is the
balanced accuracy.}
\label{fig:eval}
\end{figure}

\section{What came out}

\begin{table}[H]
\centering
\caption{Few-electron region ($n_1,n_2\le\vNmax$), \vNTest{} test devices,
means.  Full table for all 15 budgets: \texttt{few\_electron\_results.pdf}.}
\label{tab:main}
\small
\input{table_main.tex}
\end{table}

\begin{figure}[!htbp]
\centering
\includegraphics[width=\textwidth]{figures/fig_few_vs_coverage.pdf}
\caption{Every budget.  Blue: U-Net map.  Green: after the line fit.
Orange: the conventional route on the full image.  Dashed: ``all
$(0,0)$''.}
\label{fig:cov}
\end{figure}

The results, in order of importance:
\begin{enumerate}
\item \textbf{The U-Net maps alone give few correct charge states}:
  balanced accuracy \vUfourBal\,\% at the smallest budget and \vUBal\,\%
  [\vUBalLo, \vUBalHi] at the largest, and the gates land in $(1,1)$ on
  \vUfourLand--\vULand\,\% of the devices.
\item \textbf{The line fit helps every budget}, by \vGainMin{} to \vGainMax{}
  points of balanced accuracy, reaching \vFBal\,\% [\vFBalLo, \vFBalHi] at
  $8\times60$.  The intervals with and without it do not overlap, so this is
  not luck of the test set.  Landing reaches \vBestLand\,\% at best.  The
  median device still does not land in $(1,1)$.
\item \textbf{The conventional route on the full image is far better}:
  \vDBal\,\% [\vDBalLo, \vDBalHi] and \vDLand\,\% landing.  It measures
  every pixel, \vFold{} times more than the $8\times60$ U-Net.
\item \textbf{Fewer honeycombs do not rescue the U-Net}
  (Table~\ref{tab:density}).  On windows with few charge states its pixel
  accuracy is higher, but so is the ``all $(0,0)$'' score, and its state
  recall is about the same in every group.
\end{enumerate}

\begin{table}[H]
\centering
\caption{Test devices grouped by the number of charge states in the window
(whole window, $8\times60$ U-Net without the fit).}
\label{tab:density}
\small
\resizebox{\textwidth}{!}{\input{table_density.tex}}
\end{table}

\begin{figure}[!htbp]
\centering
\includegraphics[width=\textwidth]{figures/fig_fit_examples.pdf}
\caption{Three test devices: a map with gaps that the fit repairs (top), a
complete map left unchanged (middle), and a map whose lines are broken into
dashes that straight segments cannot rejoin (bottom).}
\label{fig:examples}
\end{figure}

\section{The decisive test: the conventional route at equal cost}
\label{sec:equal}

\subsection{Why this test is needed}
Comparing the U-Net with the conventional route on the \emph{full} image is
unfair to the U-Net: the full image costs \vFold{} times more measurements.
In an experiment, measurement time is exactly what one wants to save.  The
fair question, and the one Prof.\ Fujita asked, is therefore: \emph{given the
same number of measurements, does predicting the lines with the U-Net give
better charge states than the conventional route of building an image and
detecting the lines in it?}  If not, the U-Net has no reason to exist.

\subsection{What the competitors do}
For every one of the \vEqNBud{} U-Net budgets, the conventional route was
given exactly the U-Net's measurement budget, in two ways an experimentalist
might spend it (Fig.~\ref{fig:equalex}).
\begin{enumerate}
\item \textbf{Rays + interpolation} (``reconstruct, then detect''): take the
  \emph{same} ray measurements the U-Net receives, and fill the gaps between
  the rays by interpolation (nearest, linear or cubic) to get a full image.
  Then find the lines with the step detector of the full-image route, and
  read off the charge states.  This is the simplest form of reconstruction
  followed by line detection.
\item \textbf{Coarse raster scan}: spend the same number of points on a
  regular grid instead of rays.  For $8\times60$ that is \vEqNPts{} points,
  a $\vEqK\times\vEqK$ grid.  This is what one does in the lab when time is
  short.  Two variants were tried:
  \begin{itemize}
  \item upsample the coarse image to full resolution, then find the steps;
  \item find the steps on the coarse grid itself, where a transition is one
    jump between neighbouring measurements, and then map the result to full
    resolution.  This is the better way to treat a coarse scan.
  \end{itemize}
\end{enumerate}
Each competitor's line map then goes through the same labelling, with and
without the same line fit, and is scored in the same way.

\subsection{How the comparison was kept fair}
\begin{itemize}
\item \textbf{Same information.}  The ray competitor sees exactly the data
  the U-Net sees.  The coarse competitor sees the same number of points.
\item \textbf{Tuned like the U-Net.}  Each competitor's settings (the
  interpolation and the step detector's window, threshold and dilation) were
  scanned on the validation devices and the best was applied unchanged to
  the test devices, just as the U-Net's threshold and the line fit were.
  Chosen: \emph{\vEqChosen}.
\item \textbf{Same repair offered.}  The competitors were also given the
  line fit, and the better of with and without is shown for each.
\item \textbf{A bias found and corrected.}  A first run used too narrow a
  threshold range for the competitors and only the upsampled coarse
  variant.  That would have handicapped them, so it was re-run with a wider
  range and with the native coarse variant.  It made little difference: the
  best validation score of every competitor stayed between 12 and 14\,\%.
\item \textbf{A sanity check on the coarse scan.}  Even when the coarse
  route's setting was tuned on individual test devices themselves (an
  unfair advantage, used only as a check), it reached only 24\,\% and
  14\,\% on two of them.  Its weakness is real, not a tuning artefact.
\item \textbf{Paired statistics.}  The U-Net and its competitor are scored
  on the same devices, so their per-device differences are bootstrapped
  directly.  This is a sharper test than comparing two separate averages.
\end{itemize}

\subsection{Result}

\begin{table}[H]
\centering
\caption{Few-electron balanced accuracy (\%), \vNTest{} test devices, every
budget.  ``Best competitor'': the best of the six conventional variants at
that budget (chosen with hindsight on the test set, which favours the
competitor).  Difference: U-Net + fit minus best competitor, mean over
devices, with a paired 95\,\% bootstrap interval.}
\label{tab:equal}
\small
\input{table_equal.tex}
\end{table}

\begin{figure}[!htbp]
\centering
\includegraphics[width=\textwidth]{figures/equal_cost_curves.pdf}
\caption{Every method at every budget.  Green squares: U-Net + line fit.
Blue circles: U-Net alone.  Purple and orange: the conventional route with
the same number of measurements.  Orange diamond: the full image.}
\label{fig:equal}
\end{figure}

\begin{figure}[!htbp]
\centering
\includegraphics[width=\textwidth]{figures/equal_cost_examples.pdf}
\caption{The same three devices, $8\times60$ budget (\vEqNPts{}
measurements).  (a, b) U-Net + fit.  (c, d) The same rays, interpolated:
the wedges between the rays fill with smeared streaks, and the step detector
finds edges in them that are not transitions.  (e, f) A $\vEqK\times\vEqK$
coarse scan: the lines are visible but at about 5\,px per step their
positions and crossings are too coarse to count the states.  On the
dense-honeycomb device (bottom) the conventional routes do slightly better
than the U-Net.}
\label{fig:equalex}
\end{figure}

\begin{itemize}
\item \textbf{The U-Net with the line fit wins at every budget}:
  \vEqWins{} of \vEqNBud{} paired intervals lie entirely above zero.  At
  $8\times60$ it reaches \vFBal\,\% against \vEqCompBal\,\% for the best
  competitor (\vEqComp), a difference of \vEqDiff{} points
  [\vEqDiffLo, \vEqDiffHi].
\item \textbf{The competitors barely beat ``all $(0,0)$''.}  At $8\times60$:
  rays interpolated \vEqRaysBal\,\%, coarse upsampled \vEqUpBal\,\%, coarse
  native \vEqNatBal\,\%, against \vBaseBal\,\% for labelling everything
  $(0,0)$.  With \vEqNPts{} points, neither interpolation nor a coarse grid
  preserves the transitions well enough to count the states.
\item \textbf{The line fit is part of the advantage.}  Without it, the U-Net
  alone beats the best competitor's average at \vEqRawWins{} of
  \vEqNBud{} budgets.  With it, at all of them.  On the competitors the line
  fit helps little or even hurts, because their line maps are cluttered
  rather than broken.
\item \textbf{Not on every device.}  On windows with many charge states the
  conventional routes can do as well or better (Fig.~\ref{fig:equalex},
  bottom).
\end{itemize}

\section{What I understood}

\paragraph{1. The line-level score does not tell whether a map is usable.}
This is the most important observation.  Figure~\ref{fig:f1} plots, for
every test device, the U-Net map's F1@1 against the charge-state accuracy
obtained from it.  The two are only weakly related (Spearman
$\rho=\vRhoU$).  Of the \vNHiFone{} devices with F1@1 $\ge 0.85$, which
looks like an excellent line map, the balanced accuracy ranges from
\vHiFoneBalMin\,\% to \vHiFoneBalMax\,\%.  F1 counts \emph{how much} of each
line is found.  Charge states need every line to be \emph{closed}.  A map
can find 95\,\% of every line and still be useless if the missing 5\,\% are
gaps in the wrong places.  So the U-Net is trained and selected for a target
that is not the one that matters downstream.

\begin{figure}[!htbp]
\centering
\includegraphics[width=0.55\textwidth]{figures/f1_vs_states.pdf}
\caption{Each point is one test device ($8\times60$).  The line-level F1@1
of the U-Net map barely predicts how many charge states can be read off it
(blue), even after the line fit (green).}
\label{fig:f1}
\end{figure}

\paragraph{2. The labelling is not the bottleneck; the topology of the line
map is.}  On complete line maps the labelling is essentially perfect.
Every error traces back to a line that is missing, broken or drawn where
there is none.

\paragraph{3. Most errors are repairable with geometry, but not all.}
Straight-line continuation closes many gaps (doubling the balanced
accuracy), which shows the information is largely present in the U-Net
output.  Its limits are visible too: dashed lines are not rejoined, and an
occasional continuation draws a false line.  The physics offers stronger
constraints that this repair does not yet use.  In one device, all dot-1
lines share one slope, all dot-2 lines share another, and they repeat with a
regular spacing.

\paragraph{4. The U-Net's value is measurement cost, and it is real.}  On
noise-free images, finding steps in the full image recovers almost all
charge states, so the U-Net cannot win on accuracy at full information.
At \emph{equal} information it wins clearly (Sec.~\ref{sec:equal}).  The
reason is visible in Fig.~\ref{fig:equalex}.  Interpolating between rays
invents structure that is not there, and a coarse grid blurs the crossings.
The U-Net has learned what a honeycomb looks like, so it can draw the lines
\emph{between} the rays where no measurement exists.  This is the
answer to the first half of Prof.\ Fujita's question: what is gained by
predicting the lines directly is that the prior knowledge of the line
geometry fills in what was not measured.

\paragraph{5. Pixel accuracy is the wrong metric for this task.}  It
rewards the large $(0,0)$ region.  Balanced accuracy and $(1,1)$ landing
are the numbers that match what an experiment needs.

\paragraph{6. A useful by-product: the map can check itself.}  The
consistency check needs no ground truth: 1.00 on true maps, much lower on
broken ones.  In an experiment it could decide when a map is not to be
trusted, and where to measure more.

\section{Is this a scientific contribution?}
\label{sec:contribution}

\textbf{Now there is a result that can become one.}  The equal-cost test
gives a defensible statement: \emph{with the same few percent of the gate
plane, predicting the transition lines with a U-Net and repairing their
gaps identifies the few-electron charge states far better than
reconstructing an image from the same measurements and detecting lines in
it} (\vFBal\,\% against at most \vEqCompBal\,\% at $8\times60$; better at all
\vEqNBud{} budgets, with paired intervals excluding zero).  That answers
Prof.\ Fujita's question in the U-Net's favour, under the conditions
tested.

\textbf{What it is not.}  It is not a claim that the U-Net is accurate
enough for tuning.  It identifies about two thirds of the few-electron
states and lands in $(1,1)$ in fewer than half of the devices.  Measuring the
whole image is still far better.  The contribution is about \emph{cost}:
how much charge-state information survives a 20--60-fold cut in
measurements.

\textbf{What still has to be done before it is a publishable claim:}
\begin{itemize}
\item \textbf{A stronger reconstruction competitor.}  The competitors here
  are simple (interpolation, coarse grid).  The earlier probabilistic
  reconstruction of this project, or a published compressed-sensing or
  inpainting method, must be run through the same pipeline.  If a good
  reconstruction closes the gap, the conclusion changes.
\item \textbf{Noise.}  Every result here is noise-free.  Noise will hurt
  the step detector and the U-Net differently.
\item \textbf{The line fit must be part of the method.}  The advantage
  depends on it, so it has to be described and justified as part of the
  method, not as a tweak.
\item \textbf{More than one sweep, and ideally one experimental CSD.}
\item \textbf{A literature check} of existing automatic charge-state
  identification methods.  It has not been done for this report, and novelty
  must not be claimed before it.
\end{itemize}

\textbf{Other findings that are contributions in their own right:}
\begin{itemize}
\item a task-level evaluation of sparse-measurement methods by what an
  experiment needs (charge states, $(1,1)$ landing), with the finding that
  line-level F1 does not predict it (Fig.~\ref{fig:f1});
\item a label-free self-check (consistency) that says when a predicted map
  should not be trusted.
\end{itemize}

\section{What to do next, in order}
\begin{enumerate}
\item \textbf{Stronger reconstruction competitor} (no training if the
  earlier reconstruction model exists): run the project's own probabilistic
  reconstruction, at the same budgets, through the same line detection and
  labelling.
\item \textbf{Physics-constrained line fit} (no training): one shared slope
  per line family and regular spacing.
\item \textbf{Noise}: add realistic sensor noise to both routes.
\item \textbf{Train for closed lines} (needs retraining): a topology-aware
  loss, or predicting the charge states directly as a second output.
\item \textbf{Realistic capacitance range} (needs retraining): few-electron
  windows only.
\end{enumerate}

\section{Limitations of this report}
Simulation only, noise-free; one sweep; \vNTest{} test devices (intervals of
several points); the validation devices chose every setting and are not
held out; the top-right corner is assumed to be $(0,0)$, which holds in this
pool.

\end{document}
